\documentclass[11pt]{article}
\usepackage{mathblatt}
% Präambel des Setzers (werkzeuge/setzer-praeambel.tex, wortgleich); dazu \gteil (Muster T2U)
\makeatletter
\setlength{\mbpfbe}{0mm}% keine Punkte (pruefung.md § 5: nur beim Prüfstein)
\newlength{\gEin}\setlength{\gEin}{0pt}
\let\gSchrift\relax
\renewcommand{\pfkreuzzeile}[1]{\par\noindent\hangindent3.2em\hangafter1\hspace*{1.6em}$\square$\ #1\par}
\newcommand{\gkreuz}[1]{$\square$\ #1\hspace{2em}}
\newcommand{\gkreuzzeile}[2]{\par\noindent\hangindent2.2em\hangafter1\mbox{#1\,$\square$}\ #2\par}% Buchstabe, Option: eine je Zeile (neben Skizze)
% links, was man ansieht, rechts, was man tut; Spaltengrenze überall an derselben Stelle (gemeinsam.md „Satz“)
% v0.17: Fußeintrag „Nr: Wert“ – kurz in einer Box (nie getrennt), lang umbrechbar mit der Nummer am Anfang
% (vorher \mbox: lange Lösungen liefen über den Rand, AS40 138 pt); robust, weil \fusshilfe in die .aux schreibt
\NewDocumentCommand{\gfuss}{m m}{\sbox0{#1:~#2}\ifdim\wd0<0.45\textwidth\mbox{#1:~#2}\else\mbox{#1:}\nobreak\ #2\fi}
\newsavebox{\gbox}\newlength{\gLinks}
% v0.10 (Lehrer 13.10.): keine starre Halbierung – die linke Spalte ist so breit wie die Figur (+ 6 mm), die
% Antwort steht direkt daneben; ist die Figur breiter als 0,62 Zeile, steht die Antwort darunter
\newcommand{\gzwei}[2]{\par\smallskip\noindent\sbox{\gbox}{#1}%
  \setlength{\gLinks}{\dimexpr\wd\gbox+6mm\relax}%
  \ifdim\gLinks>0.62\linewidth
    \usebox{\gbox}\par\smallskip\noindent\begin{minipage}{\linewidth}#2\end{minipage}%
  \else
    \begin{minipage}[t]{\gLinks}\vspace{0pt}\usebox{\gbox}\end{minipage}%
    \begin{minipage}[t]{\dimexpr\linewidth-\gLinks\relax}\vspace{0pt}\raggedright #2\end{minipage}%
  \fi\par}
% v0.9 (12.10., pruefung.md § 5 „Marke mit B, Bedingungen rechts“): drittes Argument von pfaufg = Bedingungen
% der Prüfung (v0.10: \pfohnetr = TR durchgestrichen, „FOR“), klein grau im rechten Rand auf der Grundlinie der Nummer, Zeilen mit \\
\newcommand{\gbedrand}[2]{\ifblank{#2}{}{\rlap{\hspace*{\dimexpr#1+2mm\relax}%
  \parbox[t]{15mm}{\raggedright\scriptsize\leavevmode\color{mbgrau}\strut#2}}}}
% Teil eines Bündels (a, b …): Marke links im Rand, Bedingungen rechts im Rand, auf der Zeile des Teils
\newcommand{\gteilrand}[2]{\llap{\mbpfrandmarke{#1}\hspace*{\mbpfnr}}\gbedrand{\linewidth}{#2}}
% Teil eines Bündels: Buchstabe in eigener Spalte, Inhalt eingerückt (\gEin = Einzug, damit die Spaltengrenze von
% \gzwei an derselben Stelle bleibt)
\newenvironment{gbteil}[3]{\par\smallskip\noindent\gteilrand{#2}{#3}\makebox[\mbpfteil][l]{\textbf{#1}}%
  \begin{minipage}[t]{\dimexpr\linewidth-\mbpfteil\relax}\setlength{\gEin}{\mbpfteil}\raggedright\strut\ignorespaces}%
 {\end{minipage}\par}
% Grafik höchstens so breit wie die Spalte (Bündel nebeneinander, Koordinatensystem ohne Karo)
\newsavebox{\gpbox}
\newcommand{\gpassend}[1]{\sbox{\gpbox}{#1}\ifdim\wd\gpbox>\linewidth\typeout{GPASSEND: verkleinert \the\wd\gpbox\space auf \the\linewidth}\resizebox{\linewidth}{!}{\usebox{\gpbox}}\else\usebox{\gpbox}\fi}
% Skizzenfeld (v0.10, Lehrer 13.10.): Karo ab dem linken Rand über die Zeile, „Skizze“ klein grau oben links im Feld
% (wie \pfskizzenfeld der Vorlage); der rechte Rand bleibt den Bedingungen
\newcommand{\gskizze}[1]{\par\smallskip\noindent\begin{tikzpicture}
  \clip (0,0) rectangle ({\linewidth-1mm},{#1*10mm});
  \draw[black!16,line width=0.3pt,step=5mm] (0,0) grid ({\linewidth},{#1*10mm});
  \node[anchor=north west,inner sep=2pt,font=\scriptsize,text=mbgrau,fill=white] at (0,{#1*10mm}) {Skizze};
  \end{tikzpicture}\par}
\RenewDocumentEnvironment{pfaufg}{m m m}{%
  \begin{lrbox}{\mbaufgabebox}\begin{minipage}[t]{\linewidth}%
  \gSchrift\noindent\hspace*{\gEin}\mbpfrandmarke{#1}\gbedrand{\linewidth-\gEin}{#3}%
  \begin{minipage}[t]{\mbpfnr}\strut\textbf{#2}\end{minipage}%
  \begin{minipage}[t]{\dimexpr\linewidth-\gEin-\mbpfnr\relax}\raggedright\strut\ignorespaces}%
 {\end{minipage}%
  \end{minipage}\end{lrbox}%
  \setlength{\mbaufgabehoehe}{\dimexpr\ht\mbaufgabebox+\dp\mbaufgabebox\relax}%
  \ifdefempty{\gkopfnach}{\par\addvspace{5pt}\Needspace*{\mbaufgabehoehe}}{%
    % Stufenkopf, der mit dieser (ersten) Aufgabe zusammenbleibt: Platz für beide verlangen, dann den Kopf setzen
    % (12.10., Quadratische: ein Kopf stand allein unten, die Aufgabe mit Koordinatensystem auf der nächsten Seite)
    % v0.17 (AS40 „Terme deuten“ allein unten): hinter dem Kopf steht das Sprungziel (\special, ein Knoten, nach
    % dem die Seite brechen durfte) – darum \nobreak nach dem Kopf; Vorhalt 56 pt = Kopf mit Linie und Abständen
    \par\Needspace*{\dimexpr\mbaufgabehoehe+56pt\relax}\gkopfnach\par\nobreak\gdef\gkopfnach{}\addvspace{5pt}}%
  \noindent\usebox{\mbaufgabebox}\par\addvspace{5pt}}
% Stufenkopf vormerken: \gkopfvor{<Kopf>} – die nächste pfaufg setzt ihn, sobald beide zusammen Platz haben
\def\gkopfnach{}
\newcommand{\gkopfvor}[1]{\gdef\gkopfnach{#1}}
% Teilaufgabe mit Skizze links in gzwei: Buchstabe in eigener Spalte, darunter Text und Skizze
\newcommand{\gteil}[2]{\makebox[\mbpfteil][l]{\textbf{#1}}\parbox[t]{\dimexpr0.5\textwidth-\gEin-\mbpfnr-5mm-\mbpfteil\relax}{\raggedright #2}}
% Stufenkopf der Vorlage (mathblatt.sty 2026-10-12: Linie, Abstand, 13 pt, \pfjahre) ohne ihren Vorhalt von
% 14 Zeilen: der Setzer fordert vor dem Kopf 0,16 Seitenhöhe (Kopf + erste Aufgabe), mehr nicht
\setlength{\mbpfstufevorhalt}{0pt}
\makeatother
\newcommand{\karo}[1]{\par\smallskip\noindent\begin{tikzpicture}
  \clip (0,0) rectangle ({\linewidth-1mm},{#1*10mm});
  \draw[black!16,line width=0.3pt,step=5mm] (0,0) grid ({\linewidth},{#1*10mm});
  \end{tikzpicture}\par}
\newcommand{\kk}[1]{$\square$\,#1\hspace{0.9em}}
\newcommand{\linie}[1]{\,{\color{black!45}\rule[-1pt]{#1}{0.4pt}}\,}
\newcommand{\pkt}{\textperiodcentered{}}
% v0.14 (14.10.): Stufenkopf mit fünf Punkten (Option = Jahre für die Bestellkarte) oder „zuletzt 19“
\NewDocumentCommand{\gpunkte}{o m}{#2}
\newcommand{\gzuletzt}[1]{zuletzt #1}
% v0.7 (12.10.): Schülerfrage unter dem Stufenkopf; Ankreuzoption als Bild (Buchstabe und Kästchen darunter)
\newcommand{\gfrage}[1]{\par\nobreak\noindent\hspace*{\mbpfnr}{\small\itshape #1}\par\nobreak\addvspace{2pt}}
\newlength{\gbildbreite}\setlength{\gbildbreite}{0.32\linewidth}% Setzer setzt sie je Reihe (vier Bilder: 0,24)
\newcommand{\gbildoption}[3]{\begin{minipage}[t][#3][b]{\gbildbreite}\centering #2\par\smallskip\mbox{#1\,$\square$}\end{minipage}}
% Lineare (12.10.): Bild mit Buchstabe ohne Kästchen (Diagramm zuordnen); Aussage mit „wahr / falsch“ zum Ankreuzen
\newcommand{\gbildname}[3]{\begin{minipage}[t][#3][b]{\gbildbreite}\centering #2\par\smallskip\mbox{\textbf{#1}}\end{minipage}}
\newcommand{\gwahrfalsch}[2]{\par\noindent\hangindent2.2em\hangafter1\makebox[2.2em][l]{(#1)}#2\hfill\mbox{wahr $\square$\quad falsch $\square$}\par}
\begin{document}

\pfheftstil
\pflaufkopf[Abitur GK \pkt{} Prüfung]{Funktion aufstellen}{AA110}\pfseitevon
\pfheftkopf{Funktion aufstellen}{}
\pfstufenzeile{Funktionsgleichung gesucht \pkt{} Zielfunktion aufstellen}

% ===================== Stufe: Funktionsgleichung gesucht =====================
\gkopfvor{\pfstufekopf{Funktionsgleichung gesucht}{\gpunkte[23 22]{\textcolor{black!25}{$\circ$}\,\textcolor{black!25}{$\circ$}\,\textcolor{black!25}{$\circ$}\,\textcolor{black!80}{$\bullet$}\,\textcolor{black!80}{$\bullet$}}}\special{pdf:dest (stufe-funktionsgleichung-gesucht) [@thispage /XYZ @xpos @ypos null]}}
% kein Einstieg: keine Schülerfrage in der Gliederung (Schülerfrage, -Fälle, -Form, -Weg, -Skizze; oder Frage) – Einstieg nur von dort, nie aus der Bank (pruefung.md § 5, 12.10.)
% 2022-bebb-gk-B2.2g-m1 (Original 2022-bebb-gk-B2.2g, Schwere I, rang 1; herausgelöste Fassung: keine)
\begin{pfaufg}{GK~’22\\B2.2g}{1.}{}
Die Abbildung zeigt den Querschnitt eines Betonteils für den Hochwasserschutz. Der untere Rand ist die Strecke zwischen $C(0{,}4 \mid -0{,}8)$ und $D(4 \mid 1)$ auf der Geraden $g$. Weise nach, dass $g(x) = \tfrac{1}{2}x - 1$ eine Gleichung von $g$ ist.
\gzwei{\begin{abiksys}[xmin=-0.5,xmax=4.5,ymin=-1.5,ymax=2.5,le=1.2,ystep=0.5]\funktionab[2.6]{-\x^3/16+3*\x^2/8}{h}{0}{4}\draw[line width=0.8pt] (0.4,-0.8) -- (4,1) node[midway,below right] {$g$};\draw[line width=0.6pt] (0,0) -- (0.4,-0.8) (4,2) -- (4,1);\punkt{0}{0}{A}\punkt{4}{2}{B}\punkt{0.4}{-0.8}{C}\punkt{4}{1}{D}\end{abiksys}}{\karo{3}}
\fusshilfe{\gfuss{1}{$\tfrac{1}{2}$; $\tfrac{1}{2}$ x $-$ 1}}
\end{pfaufg}
% 2023-bebb-gk-B2.2k-m1 (Original 2023-bebb-gk-B2.2k, Schwere II, rang 2; herausgelöste Fassung: keine)
\begin{pfaufg}{GK~’23\\B2.2k}{2.}{}
Für eine neue Backofenheizung wird der Temperaturverlauf in der Startphase durch eine ganzrationale Funktion $k$ dritten Grades beschrieben ($x$ in Minuten, $k(x)$ in °C). Es gilt: I $k(0) = 20$, II $k'(0) = 120$, III $H(6 \mid 272)$ ist Hochpunkt des Graphen von $k$. Ermittle eine Funktionsgleichung von $k$.
\karo{3}
\fusshilfe{\gfuss{2}{120; $-$120; x³ $-$ 19x² + 120x + 20}}
\end{pfaufg}
% 2022-bebb-gk-B2.2k-m1 (Original 2022-bebb-gk-B2.2k, Schwere III, rang 3; herausgelöste Fassung: keine)
\begin{pfaufg}{GK~’22\\B2.2k}{3.}{}
Der obere Rand eines größeren Betonteils verläuft zwischen $R(2 \mid 1)$ und $S(4 \mid 3{,}5)$ geradlinig; links und rechts schließen sich Kurvenstücke an, die knickfrei in $(0 \mid 0)$, $R$, $S$ und $(6 \mid 4{,}5)$ übergehen (Abbildung). Das Kurvenstück von $(0 \mid 0)$ bis $R$ soll Graph einer Funktion $k$ mit $k(x) = a \cdot x^4 + b \cdot x^2$ sein. Weise nach, dass $k$ dafür verwendet werden kann, und bestimme $a$ und $b$.
\gzwei{\begin{abiksys}[xmin=0,xmax=6.5,ymin=0,ymax=5,le=0.8]\funktionab{\x^4/64+3*\x^2/16}{}{0}{2}\draw[mbterm,line width=0.8pt] (2,1) -- (4,3.5);\funktionab{4.5-(6-\x)^4/64-3*(6-\x)^2/16}{}{4}{6}\punkt{2}{1}{R}\punkt{4}{3.5}{S}\end{abiksys}}{\karo{3}}
\fusshilfe{\gfuss{3}{0 gelten für jede Wahl; 1,25; $\tfrac{3}{16}$}}
\end{pfaufg}

% ===================== Stufe: Zielfunktion aufstellen =====================
\gkopfvor{\pfstufekopf{Zielfunktion aufstellen}{\gpunkte[24 22]{\textcolor{black!25}{$\circ$}\,\textcolor{black!25}{$\circ$}\,\textcolor{black!80}{$\bullet$}\,\textcolor{black!25}{$\circ$}\,\textcolor{black!80}{$\bullet$}}}\special{pdf:dest (stufe-zielfunktion-aufstellen) [@thispage /XYZ @xpos @ypos null]}}
% kein Einstieg: keine Schülerfrage in der Gliederung (Schülerfrage, -Fälle, -Form, -Weg, -Skizze; oder Frage) – Einstieg nur von dort, nie aus der Bank (pruefung.md § 5, 12.10.)
% 2024-bebb-gk-B2.2g-m1 (Original 2024-bebb-gk-B2.2g, Schwere II, rang 4; herausgelöste Fassung: keine)
\begin{pfaufg}{GK~’24\\B2.2g}{4.}{}
Gegeben ist die Funktion $f$ mit $f(x) = 0{,}5x^4 - 3x^2 + 2{,}5$, $x \in \mathbb{R}$. Für $-1 \le x \le 1$ wird $f$ durch $p$ mit $p(x) = -2{,}5x^2 + 2{,}5$ angenähert. Es gibt genau zwei Stellen in $[-1; 1]$, an denen die Differenz von $p$ und $f$ am größten ist. Zeige, dass diese größte Differenz $\tfrac{1}{8}$ beträgt.
\karo{3}
\fusshilfe{\gfuss{4}{0,5x² $-$ 0,5x$^{4}$; $\pm$1/$\sqrt{2}$; größte Differenz $\tfrac{1}{8}$}}
\end{pfaufg}
% 2022-bebb-gk-B2.1h-m1 (Original 2022-bebb-gk-B2.1h, Schwere II, rang 5; herausgelöste Fassung: keine)
\begin{pfaufg}{GK~’22\\B2.1h}{5.}{}
Gegeben ist die in $\mathbb{R}$ definierte Funktion $f$ mit $f(x) = (x + 3) \cdot e^{-x}$ und $f'(x) = -(x + 2) \cdot e^{-x}$. Der Ursprung und der Punkt $P(u \mid v)$ des Graphen mit $u > 0$ sind gegenüberliegende Ecken eines Rechtecks, dessen Seiten parallel zu den Koordinatenachsen sind. Genau für ein $u$ ist der Flächeninhalt maximal. Zur Kontrolle: $A'(u) = (3 - u - u^2) \cdot e^{-u}$. Ermittle die Koordinaten von $P$ für den größten Flächeninhalt.
\karo{3}
\fusshilfe{\gfuss{5}{(u² + 3u) · $e^{-u}$; 1,30; P(1,30|1,17)}}
\end{pfaufg}

\pfblattende{10}{26 & 25 & 24 & 23 & 22 & 21 & 20 & 19 & 18 & 17}{ &  & B2.2g & B2.2k & B2.1h B2.2g B2.2k & A1.3b B2.1k &  & B2.1c & B1.1e B1.1g & B1.1f}%
{\pfausschnitt{Wendepunkte \pkt{} Gleichungen und Ungleichungen}{Funktion aufstellen}{Verschieben, spiegeln, strecken \pkt{} Graph im Sachzusammenhang}}
\end{document}